\section{Proof details}

\subsection{Normal convergence and differentiation}

Fix a compact set \(K\subset\{s:\Re s>\sigma_*\}\), and let
\(\sigma_0=\min_{s\in K}\Re s\).  Equation (13) gives
\[
 \sum_{n\ge1}\sup_{s\in K}
 \frac{|\Tr\LL_s^n|}{n}
 \le
 \sum_{n\ge1}
 \frac{(2\alpha_0^{\sigma_0})^n}
 {n(1-\alpha_0^n)}<\infty.
\]
Thus the trace logarithm is a normally convergent series of holomorphic
functions.  Cauchy's integral formula on a slightly larger compact
neighborhood permits termwise differentiation, which proves (17).  For real
\(\sigma\), \(0<\ee^{-T_\omega}<1\) and
\(\epsword_\omega\in\{-1,1\}\), so every denominator and every numerator in
(17) is positive.

\subsection{The determinant lower anchor}

For real \(\sigma>\sigma_*\), put
\[
 S(\sigma)=\sum_{n\ge1}\frac1n
 \sum_{\omega\in\{L,R\}^n}
 \frac{\ee^{-\sigma T_\omega}}
 {1-\epsword_\omega\ee^{-T_\omega}}.
\]
Termwise positivity and (14) give
\(0<S(\sigma)\le B(\sigma)\).  Equation (15) becomes
\(\Dpol(\sigma)=\ee^{-S(\sigma)}\), hence
\[
 \ee^{-B(\sigma)}\le\Dpol(\sigma)<1.
\]
At \(\sigma=2\), \(q_2=2\alpha_0^2\), so \(B(2)=B_2\).  Combining this
anchor with the pure-left term in (17) proves (20).

\subsection{Cauchy transfer}

For \(0<r<R-2\), Cauchy's estimate on \(|s-2|\le r\) gives
\[
 |\Dpol'(2)|\le
 \frac{\max_{|s-2|\le r}|\Dpol(s)|}{r}
 \le\frac{\Mdet(R)}{r}.
\]
Letting \(r\uparrow R-2\) proves (21).  This centered-disk argument avoids
the extraneous factor that would arise from estimating the derivative
directly on the circle \(|s|=R\).

\subsection{Super-polynomial maximum modulus}

If an entire function \(f(s)=\sum_{m\ge0}a_ms^m\) is not a polynomial, then
for every \(A>0\) there is an integer \(m>A\) with \(a_m\ne0\).  Cauchy's
coefficient estimate on \(|s|=R\) gives
\(|a_m|\le M_f(R)R^{-m}\).  Therefore
\(M_f(R)/R^A\ge|a_m|R^{m-A}\to\infty\).  Applying this observation to
\(f=\Dpol\) completes the proof of (9).

